% chapters/00_manifest.tex
% Part 0: System Manifest & Context Handoff (v5.5)

\chapter*{System Manifest \& Context Handoff (v5.5)}
\addcontentsline{toc}{chapter}{System Manifest \& Context Handoff (v5.5)}
\markboth{System Manifest}{System Manifest}

\section*{Document Identification}

\begin{description}[leftmargin=0pt,style=nextline]
\item[Title]
NRMO Complete System v5.5 --- Unified Specification
\item[Subtitle]
Mathematical Foundations, Cognitive Architecture, and Empirical Validation
\item[Author]
Takashi Ikeya (Zarame), NRMO Research Framework
\item[Version]
5.5 (April 2026 --- Empirical Validation Integrated)
\item[Predecessors]
v4 (\textit{Unified Specification v4 --- Complete English Edition},
February 2026, 222\,pp.); v2.0 (March 2025, foundational architecture);
v2.1 (May 2025, secretary modules and cognitive layer); v3.0--v4.0 (2025--2026,
incremental refinement and structural consistency).
\item[Build context]
v5.5 integrates four research strands developed in parallel between
February and April 2026: the formal arXiv submission
(\cite{Ikeya2026arXiv}), the NRMO vNext Civilisation Simulator
(\cite{Ikeya2026vNext}), the StrongEngine $\Omega$ Full integrated
specification (\cite{Ikeya2026Integrated}), and the v2.5 reference
implementation (this work, Part X).
\end{description}

\section*{What is New in v5.5}

\nrmobox{Summary of v5.5 additions}{
\begin{itemize}[leftmargin=1.2em]
\item \textbf{Part VII:} Mathematical Formalisation. Four theorems (T1
  Ruin Dilution, T2 CMDP Separation, T3 Variational Analysis, T4 CVaR
  Tractability), each proved with full rigour, establish that NRMO is
  not a heuristic but a strictly distinct decision-theoretic regime.
\item \textbf{Parts VIII--IX:} the vNext Civilisation Simulator and
  StrongEngine $\Omega$ Full are documented with complete state vectors,
  five world families, nine comparison strategies, an adaptive tuning
  layer, and a 30-civilisation casebook.
\item \textbf{Part X:} the engineering implementation v2.5 is presented
  in full, including the pipeline architecture (NRMO Kernel + Adaptive
  Tuning + MAPLayer + Shinobi Fleet + Norn/Skuld + $\Omega$ Full), the
  module structure, and the single one-line change from v2.1 that
  produced v2.5.
\item \textbf{Part XI:} empirical validation results from $n = 250$
  five-world episodes (April 2026), with statistical comparison against
  the v5.2 baseline, world-by-world standard errors, and death-cause
  analysis.
\item \textbf{Part XII:} methodological notes documenting five failed
  approaches to PlanetaryStress recovery, the Option C hybrid
  experiment via monkey-patch adapter, and the verification tools used
  to ensure measurement integrity.
\item \textbf{Appendices A--D:} full code listing, world parameter
  schemas, the 30-civilisation roster with profiles, and a session-by-session
  development log.
\end{itemize}
}

\noindent
Parts I through VI preserve the cognitive architecture finalised in v4
(the v5\_updated PDF) without substantive change. The preservation is
deliberate: the cognitive layer (Type ZERO, Passive Pattern, HST-N,
TTM/PPS, A\_allowed, the secretary modules, governance toolkit, and
APCSO) is a self-contained body of work whose validity does not depend
on the engineering implementation. Light copyediting has been applied
for cross-referencing into the new parts; no chapter has been deleted,
reordered, or rewritten beyond what is required for typographic
consistency. A complete change register appears in
Appendix~\ref{app:session-log}.

\section*{What v5.5 Is Not}

This document is neither a textbook nor a self-help system. Several
qualifications apply:

\begin{description}
\item[Not a doctrine of fear.] NRMO does not say ``minimise risk.'' It
says: \emph{do not enter states from which no future is reachable; within
the non-ruin domain, maximise long-horizon optionality}. The two
clauses are lexicographically ordered, but both are operative.
\item[Not a generic productivity framework.] The framework is designed
for environments with absorbing failure states. In environments where
ruin is bounded or recoverable, the constraint is slack and the
framework reduces to standard expected-utility optimisation.
\item[Not a complete formal decision theory.] T1--T4 establish the
ruin-avoidance regime is well-defined and tractable; they do not solve
optimal exploration in arbitrary MDPs. Open problems are catalogued in
Part XII.
\item[Not a substitute for judgement.] The cognitive architecture
disciplines decision-making but does not replace it. The secretary
modules detect distortion; they do not resolve it. The simulator
provides counterfactuals; it does not provide truth.
\item[Not validated by historical case studies.] Parts VIII--XI use
synthetic civilisation simulation. Mappings to historical civilisations
appear illustratively in Part IX (the 30-civilisation casebook) but are
not claimed as historical proof.
\end{description}

\section*{Reading Paths}

\noindent
The document is large (approximately 300 pages). Different readers
should enter at different points.

\begin{description}
\item[Mathematician / decision theorist.] Read Part VII first. The
four theorems are self-contained. Cross-reference Part VIII for the
state-space instantiation when needed. Skip Part II entirely if not
interested in the cognitive architecture.

\item[Practitioner adopting NRMO for personal operation.] Read Part 0
(this section), then Parts I and II. Part II Chapter 12 (Operational
Refinements) and Chapter 18 (NRMO Life SOP v2.0) are the most
immediately actionable. The simulator parts (VIII--XI) can be deferred.

\item[Researcher building on the simulator.] Read Part 0, then Parts
VIII--X in order. Part XI provides the empirical baseline for
extensions. Part XII documents which approaches have failed and why.
The code listing in Appendix A is the canonical implementation.

\item[Reader integrating with existing safety frameworks.] Read Part
VII Chapters 39--40 (T1 and T2). Theorem 2 establishes a strict
separation from Lagrangian CMDPs that is relevant to anyone
implementing constrained policy optimisation.

\item[Complete reading.] In order, Part 0 through Part XII. The
typical pace for a careful first reading is 8--12 hours over multiple
sessions.
\end{description}

\section*{Architectural Invariants}

The following statements hold throughout v5.5 and must not be violated
by any extension. They are restated formally in Part II Chapter 9
(Theoretical Invariants) and Part IX Chapter 50 (Engine Specification).

\begin{invariant}[Governance--Execution Separation]
\label{inv:gov-exec-sep}
The NRMO core governance layer defines admissibility:
$A_t = \mathrm{NRMO}(X_t)$. The execution layer (StrongEngine $\Omega$
Full, the engineering pipeline, or any future engine) selects an action
within admissibility: $a_t = \Omega(A_t)$ with $a_t \in A_t$ as a
hard requirement. The governance and execution layers do not share an
objective function and do not communicate via shared mutable state. The
governance layer never performs ranking; the execution layer never
modifies $A_t$.
\end{invariant}

\begin{invariant}[True Ruin as Hard Constraint]
\label{inv:true-ruin-hard}
Any candidate action that increases $\Pr(\tau^F \le H)$ above the
admissibility threshold $\varepsilon_t$ is rejected by the governance
layer regardless of its expected value. The threshold $\varepsilon_t$
is itself state-responsive (Part II Chapter 12) but never relaxed below
the configured floor $\varepsilon_{\mathrm{floor}}$ in any non-ruin
state.
\end{invariant}

\begin{invariant}[Optionality is the Long-Horizon Objective]
\label{inv:optionality}
Within the non-ruin domain, the primary maximisation target is
optionality: the preserved breadth of future action space. Survival is
a constraint; optionality is the objective. A policy that survives
without preserving optionality is in passive ruin
(Part II Chapter 7, Part VIII Chapter 47).
\end{invariant}

\begin{invariant}[Intelligence is not Authority]
\label{inv:intelligence-not-authority}
Search capacity, strategy generation, and computational resource use
are downstream of governance. An execution engine more capable than
NRMO core does not earn authority over the boundary. This invariant is
the structural reason fast-acting components ($\Omega$ Full's mutation,
synthesis, and invention layers) cannot bypass slow-acting components
(NRMO veto and passive ruin detection).
\end{invariant}

\section*{Versioning and Change Discipline}

\noindent
v5.5 is a \emph{compatible extension} of v4. No invariant has been
weakened; no chapter content has been removed; no terminology has been
redefined. The numbering convention is:

\begin{itemize}
\item Major revision (e.g., v4 $\to$ v5.0) signals a substantial
  architectural addition. v5.0 was not released; v5.5 is the first
  released revision after v4 because it integrates the formal
  mathematics, simulator, and engineering implementation into one
  document.
\item Minor revision (e.g., v5.5 $\to$ v5.6) would signal new empirical
  results, refinements to the engineering layer, or corrections to
  proofs --- without breaking compatibility.
\item Breaking change (e.g., v5.5 $\to$ v6.0) would signal modification
  of an architectural invariant. No such change is anticipated.
\end{itemize}

\noindent
The sentinel statement at the end of every revision is:

\nrmobox{Closing statement}{
\centering
\textit{It no longer says only ``do not die.''}\\
\textit{It says: ``do not close the future.''}
}

\section*{Acknowledgements}

\noindent
The cognitive architecture (Parts I--VI) was developed over the period
2024--2026 in close consultation with multiple instances of the Claude
language model (Anthropic), serving as interpretive partner and
counterfactual reasoning sandbox. The mathematical formalisation (Part
VII) was completed February 2026, with three errors in the v2 draft
corrected through external feedback. The simulator and engineering
implementation (Parts VIII--XI) were developed March--April 2026 over
eight intensive working sessions, also in collaboration with Claude.
The methodological discipline of Part XII --- particularly the
requirement to verify measurement integrity before drawing conclusions
--- is a direct result of those sessions.

\noindent
Section-level errata, addenda, and discussions are tracked in the
public NRMO Research Framework repository. The reference
implementation (\texttt{nrmo\_full v2.5}) is provided as
Appendix~\ref{app:code}.

\vspace{1em}
\hrule
\vspace{1em}
\noindent
\textit{End of System Manifest. The next chapter (Part I) begins the
preserved core from v4.}
